% Part: second-order-logic
% Chapter: metatheory
% Section: introduction

\documentclass[../../../include/open-logic-section]{subfiles}

\begin{document}

% SEGMENT OLP-0331-S01

\olfileid{sol}{met}{int}

\olsection{அறிமுகம்}

முதல்தரத் தருக்கத்திற்குப் பல நல்ல பண்புகள் உள்ளன. அது
தீர்மானிக்கத்தக்கது அன்று என்பது நமக்குத் தெரியும்; ஆனால் குறைந்தது அது
அடிகோளாக்கத்தக்கது. அதாவது, முதல்தரத் தருக்கத்திற்கு ஒலிப்பும்
நிறைவும் உடைய நிறுவல் முறைமைகள் உள்ளன. வேறு சொற்களில், எந்த
!!{sentence}s இன் கணம்~$\Gamma$ க்கும் !!{sentence}~$!A$ க்கும்,
$\Gamma \Entails !A$ எனில், எனில் மட்டுமே, $\Gamma \Proves !A$
என்ற பண்புடைய !!a{derivability} உறவு~$\Proves$ ஐ அவை அளிக்கின்றன.
குறிப்பாக, முதல்தரத் தருக்கத்தின் ஏற்புடைமைகள் !!{computably
  enumerable} என்பதை இது குறிக்கிறது. $f$ இன் மதிப்புகள் அனைத்தும்
மொழி~$\Lang{L}$ இன் ஏற்புடைய !!{sentence}s ஆகவும் அவை மட்டுமே
ஆகவும் இருக்கும் கணிக்கத்தக்க சார்பு~$f\colon \Nat \to \Sent[L]$
ஒன்று உள்ளது. ஏனெனில், !!{derivation}s ஐ எண்ணிடலாம்; ஒற்றை
!!{sentence} ஒன்றை !!{derive} செய்பவை பின்னர் அந்த !!{sentence} க்கு
அனுப்பப்படுகின்றன. இரண்டாம்தரத் தருக்கம் முதல்தரத் தருக்கத்தைவிட
அதிகமாக வெளிப்படுத்த வல்லது; எனவே அதன் ஏற்புடைமைகளைப் பொதுவாகக்
கைப்பற்றுவது மிகவும் சிக்கலானது. உண்மையில், இரண்டாம்தரத் தருக்கம்
தீர்மானிக்கவியலாதது மட்டுமன்று; அதன் ஏற்புடைமைகள் !!{computably
  enumerable} கூட அல்ல என்று காட்டுவோம். ஆகவே இரண்டாம்தரத்
தருக்கத்திற்கு ஒலிப்பும் நிறைவும் உடைய நிறுவல் முறைமை இருக்க முடியாது
(ஒலிப்புடைய, ஆனால் நிறைவற்ற நிறுவல் முறைமைகள் கிடைக்கின்றன; அவை
ஆராய்ச்சியின் முக்கியப் பொருள்களாகவும் உள்ளன).

முதல்தரத் தருக்கத்திற்கு மேலும் இரண்டு பண்புகள் உள்ளன: அது
இறுக்கமானது (!!{sentence}s இன் கணம்~$\Gamma$ இன் ஒவ்வொரு
முடிவுறு உட்கணமும் நிறைவுறுத்தத்தக்கது எனில், $\Gamma$ தானும்
நிறைவுறுத்தத்தக்கது); மேலும் அதற்கு லோவென்ஹைம்--ஸ்கோலம் தேற்றம்
பொருந்துகிறது ($\Gamma$ க்கு முடிவிலா மாதிரி இருந்தால் அதற்கு
!!a{denumerable} மாதிரியும் உண்டு). இந்த இரு முடிவுகளும் இரண்டாம்தரத்
தருக்கத்தில் தோல்வியடைகின்றன. இதற்கும் காரணம், முதல்தரத் தருக்கத்தால்
வெளிப்படுத்த முடியாத !!{domain}s இன் அளவு பற்றிய உண்மைகளை
இரண்டாம்தரத் தருக்கத்தால் வெளிப்படுத்த முடிவதே.

\end{document}
